\section*{Capítulo \ref{ch:g12:catalysis} --- A catálise}

\begin{solution}{exo:g12:catalysis:1}
\omterm{def:g7:chemical-reactions:reactant}{Reagente}: \ce{H2O2}. Produtos: \ce{H2O} e \ce{O2}. \omterm{def:g12:catalysis:catalyst}{Catalisador}: \ce{I-}
(consumido na etapa 1, devolvido na etapa 2). Intermediário: \ce{IO-}
(formado na etapa 1, consumido na etapa 2).
\end{solution}

\begin{solution}{exo:g12:catalysis:2}
Homogênea; heterogênea; heterogênea; homogênea.
\end{solution}

\begin{solution}{exo:g12:catalysis:3}
O caminho não catalisado. Os níveis inicial e final são os mesmos: os
mesmos produtos se formam, com a mesma energia liberada; só o caminho
difere.
\end{solution}

\begin{solution}{exo:g12:catalysis:4}
Ele é consumido numa etapa e devolvido em outra: aparece dos dois lados
da soma das etapas, e se cancela.
\end{solution}

\begin{solution}{exo:g12:catalysis:5}
A reação acontece na superfície: quanto maior a superfície, mais
\omterm{def:g7:atoms-and-molecules:molecule}{moléculas} podem reagir ao mesmo tempo, para a mesma massa de
\omterm{def:g12:catalysis:catalyst}{catalisador} (muitas vezes caro).
\end{solution}

\begin{solution}{exo:g12:catalysis:6}
Soma: \ce{2Fe^{3+} + 2H2O2 + 2Fe^{2+} + 2H+ -> 2Fe^{2+} + O2 + 2H+ +
2Fe^{3+} + 2H2O}. Os \omterm{def:g9:ions:ion}{íons} ferro e os \omterm{def:g9:ions:ion}{íons} hidrogênio aparecem dos dois
lados e se cancelam: \ce{2H2O2 -> 2H2O + O2}.
\end{solution}

\begin{solution}{exo:g12:catalysis:7}
A metade de \qty{72}{mL} é \qty{36}{mL}: alcançada depois de cerca de
\qty{35}{min} sem \omterm{def:g12:catalysis:catalyst}{catalisador} e de \qty{3.5}{min} com ele. O \omterm{def:g12:catalysis:catalyst}{catalisador}
encurta esse tempo dez vezes.
\end{solution}

\begin{solution}{exo:g12:catalysis:8}
A reação acontece na superfície da platina; o chumbo a recobre e fica
ali, de modo que os gases não alcançam mais o \omterm{def:g9:periodic-table-first-look:metal}{metal}. O \omterm{def:g12:catalysis:catalyst}{catalisador} está
``envenenado'': presente, mas inútil.
\end{solution}

\begin{solution}{exo:g12:catalysis:9}
Até cerca de \qty{37}{\celsius}, a reação acelera com a temperatura,
como qualquer reação. Acima disso, a \omterm{def:g12:catalysis:enzyme}{enzima}, uma proteína, perde a sua
forma e, com ela, a sua atividade; a \qty{80}{\celsius}, ela está
destruída.
\end{solution}

\begin{solution}{exo:g12:catalysis:10}
\ce{2CO + 2NO -> 2CO2 + N2}: C 2 = 2, O 4 = 4, N 2 = 2. O monóxido de
carbono é oxidado pelo monóxido de nitrogênio, que é reduzido: cada
poluente elimina o outro.
\end{solution}

\begin{solution}{exo:g12:catalysis:11}
$n(\ce{O2}) = 0.072 / 24.1 = \qty{2.99e-3}{mol}$;
$n(\ce{H2O2}) = 2 \times \num{2.99e-3} = \qty{5.98e-3}{mol}$ em
\qty{10.0}{mL}: $c = \qty{0.598}{mol/L}$.
\end{solution}

\begin{solution}{exo:g12:catalysis:12}
\ce{C2H5OH -> C2H4 + H2O}, uma \omterm{def:g12:curly-arrows:categories}{eliminação}; \ce{C2H5OH -> CH3CHO + H2}.
Um \omterm{def:g12:catalysis:selective}{catalisador seletivo} acelera uma de várias reações possíveis muito
mais que as outras, e assim decide o produto.
\end{solution}

\begin{solution}{exo:g12:catalysis:13}
O \omterm{def:g12:catalysis:catalyst}{catalisador} não muda o \omterm{def:g10:reaction-progress-table:system}{estado final}: as curvas com e sem \omterm{def:g12:catalysis:catalyst}{catalisador}
terminam no mesmo volume de gás oxigênio. Obtém-se a mesma quantidade de
produto, só que mais cedo.
\end{solution}

\begin{solution}{exo:g12:catalysis:14}
$\num{6.02e23} / \num{5e6} = \qty{1.2e17}{s}$, uns quatro bilhões de
anos. Para fazer isso em um segundo: \num{1.2e17} \omterm{def:g7:atoms-and-molecules:molecule}{moléculas} de \omterm{def:g12:catalysis:enzyme}{enzima},
isto é, \qty{2e-7}{mol}.
\end{solution}

\begin{solution}{exo:g12:catalysis:15}
Amônia \ce{NH3}: $150 \times 17.0 / 14.0 = \qty{182}{Mt}$. \omterm{def:g7:atoms-and-molecules:atom}{Átomos} de
nitrogênio: $\num{1.50e14} / 14.0 = \qty{1.07e13}{mol}$, isto é,
\qty{5.36e12}{mol} de \ce{N2}.
\end{solution}

\begin{solution}{pb:g12:catalysis:1}
\textbf{1.} \ce{2H2O2 -> 2H2O + O2}.

\textbf{2.} Sim: o peróxido de hidrogênio é o \omterm{def:g11:redox:oxidant}{oxidante} do par
\ce{H2O2}/\ce{H2O} e o \omterm{def:g11:redox:oxidant}{redutor} do par \ce{O2}/\ce{H2O2}; ele oxida e
reduz a si mesmo.

\textbf{3.} Sem \omterm{def:g12:catalysis:catalyst}{catalisador}, a reação é extremamente lenta.

\textbf{4.} $20 / 22.4 = \qty{0.893}{mol}$.

\textbf{5.} Dois \omterm{def:g10:the-mole:amount}{mols} de peróxido de hidrogênio por \omterm{def:g10:the-mole:amount}{mol} de gás
oxigênio: \qty{1.79}{mol}.

\textbf{6.} \qty{1.79}{mol/L}.

\textbf{7.} $M = 2 \times 1.0 + 2 \times 16.0 = \qty{34.0}{g/mol}$;
$C_m = 1.79 \times 34.0 = \qty{60.7}{g/L}$.

\textbf{8.} A metade: \qty{0.893}{mol/L}.

\textbf{9.} Catalase: uma \omterm{def:g12:catalysis:enzyme}{enzima}, dissolvida, um \omterm{def:g12:catalysis:catalyst}{catalisador} biológico;
\omterm{def:g9:ions:ion}{íons} ferro(III): \omterm{def:g12:catalysis:homogeneous}{catálise homogênea}.

\textbf{10.} \ce{2Fe^{3+} + H2O2 -> 2Fe^{2+} + O2 + 2H+} e
\ce{2Fe^{2+} + H2O2 + 2H+ -> 2Fe^{3+} + 2H2O}; a sua soma é
\ce{2H2O2 -> 2H2O + O2}.

\textbf{11.} Com um \omterm{def:g12:catalysis:catalyst}{catalisador}, o volume sobe muito mais depressa; as
duas curvas se estabilizam no mesmo volume final.

\textbf{12.} Os \omterm{def:g9:ions:ion}{íons} ferro(III) (laranja) são transformados em ferro(II)
(verde-claro) na primeira etapa e de novo em ferro(III) na segunda; no
final, todo o ferro voltou a ser ferro(III).

\textbf{13.} A fervura destruiu a \omterm{def:g12:catalysis:enzyme}{enzima}.

\textbf{14.} $0.100 \times 20 = \qty{2.0}{L}$.

\textbf{15.} $n(\ce{O2}) = 0.100 \times 0.893 = \qty{0.0893}{mol}$, isto
é, $0.0893 \times 24.1 = \qty{2.15}{L}$.

\textbf{16.} A decomposição lenta libera gás oxigênio, que aumentaria a
pressão num frasco bem fechado.

\textbf{17.} $1.50 / 2 = \qty{0.750}{mol}$ de gás oxigênio por litro,
isto é, $0.750 \times 22.4 = \qty{16.8}{L}$: ``16.8 volumes''.

\textbf{18.} $(1.79 - 1.50)/1.79 \approx \qty{16}{\percent}$.

\textbf{19.} \qty{1.79}{mol/L}.
\end{solution}
